\documentclass[12pt,a4paper]{article}
\usepackage[utf8]{inputenc}
\usepackage[T1]{fontenc}
\usepackage[spanish]{babel}
\usepackage{geometry}
\usepackage{xcolor}
\usepackage{enumitem}

\geometry{a4paper, margin=1in}

\title{Reporte de Revisión de Pull Request \\ \large Módulo: Escaneo de Puertos (Grupo 4)}
\author{Prof. César Rodríguez \\ \small Curso: Seguridad Informática}
\date{\today}

\begin{document}

\maketitle

\section*{Estado de la Integración: \textcolor{green!70!black}{Aprobado e Integrado}}

\noindent \textbf{Calificación Final:} \textbf{9 en base a 10}

\vspace{0.5cm}
Estimados estudiantes del Grupo 4:

Se ha auditado su envío para el módulo de Escaneo de Puertos (\texttt{scanning.py}). Me complace informarles que el código ha sido \textbf{aprobado y fusionado} en la rama principal del proyecto.

\section*{Aspectos Destacados (Puntos Positivos)}
\begin{itemize}[label=$\bullet$]
    \item \textbf{Implementación Real:} Lograron reemplazar exitosamente la simulación por un escaneo de puertos TCP real utilizando la librería nativa \texttt{socket}.
    \item \textbf{Manejo Eficiente:} El uso de \texttt{s.connect\_ex()} para capturar el código de estado sin levantar excepciones por cada puerto cerrado demuestra un excelente nivel técnico.
    \item \textbf{Resiliencia:} La configuración de un \texttt{timeout} (1.0s) asegura que el orquestador no se quede colgado esperando respuestas de hosts inactivos.
    \item \textbf{Manejo de Errores:} Atrapan correctamente cualquier fallo general en el bloque \texttt{try-except} principal, resguardando la ejecución de la suite.
\end{itemize}

\section*{Motivo de la Penalización (Corrección Menor)}
El envío fue casi perfecto, pero incurrieron en una infracción leve al \textbf{Contrato de Datos (\texttt{schema\_resultados.json})}:
\begin{itemize}[label=$\bullet$]
    \item En la clave \texttt{'estudiante'}, utilizaron el nombre completo (\texttt{'Andrés Maldonado'}) en lugar del identificador estándar requerido (\texttt{'E1'}).
\end{itemize}

Dado que el resto del módulo tiene una calidad excepcional, la administración aplicó esta mínima corrección directamente en el servidor antes de la integración para no retrasar el proyecto.

\vspace{1cm}
\noindent ¡Excelente trabajo! El módulo de Escaneo de Puertos se da por completado y la suite ya es capaz de interactuar con objetivos reales en la red.

\end{document}